% GENERATED FILE. Edit canonical lessons, then run npm run generate:books.
% canonical-source-hash: fnv1a64:dc90a97172133e12
% canonical-lessons: MW-C68-jangal, MW-C68-nadi, MW-C68-bhato, MW-C68-ret, MW-C68-khet

\chapter{Forest, River, Stone, Sand, Field}
\label{ch:mw-forest-river-stone-sand-field}
\hlchaptermodality{68}

\begin{chapteropening}
\textbf{By the end of this chapter:} \emph{I can say a forest, a river, a stone, sand and a field.}

Complete the last lesson of chapter 68: I can say a forest, a river, a stone, sand and a field.
\end{chapteropening}

\section[\texorpdfstring{\mw{जंगळ} (jaṅgaḷ)}{जंगळ (jaṅgaḷ)}]{\mw{जंगळ} (jaṅgaḷ) — a forest}

\label{lesson:MW-C68-jangal}



\begin{quote}
Before the new one: say the Marwadi for a leaf, then the Marwadi for a hill.
\end{quote}

\subsection*{You'll want to know: \mw{जंगळ}}
\textbf{\mw{जंगळ}} — \emph{jaṅgaḷ} — \textquotedblleft{}a forest\textquotedblright{}.

\begin{cousinweb}
\textbf{\mw{जंगळ}} (\emph{jaṅgaḷ}) is a tadbhava, a word inherited from Sanskrit and reshaped by sound change. It comes from \emph{jāṅgala}, \textquotedblleft{}arid, sparingly grown with trees\textquotedblright{}. CDIAL gives no Marwari form in this entry (CDIAL 5177); the match is through Hindi \emph{jaṅgal}. CDIAL notes evidence for a non-Aryan origin. The sense moved from \textquotedblleft{}dry, open country\textquotedblright{} to \textquotedblleft{}forest\textquotedblright{}.
\end{cousinweb}

\begin{grammarlens}[title={What you've built}]
One more word for this chapter.
\end{grammarlens}

\subsection*{Your turn}
\begin{itemize}
\raggedright
  \item \emph{Say it:} \emph{jaṅgaḷ}
  \item \emph{Say it:} \emph{jaṅgaḷ}, once more
  \item \emph{Say it:} say \emph{jaṅgaḷ}
  \item \emph{Recall:} say the Marwadi for a leaf, then the Marwadi for a hill, then say \emph{jaṅgaḷ} again
\end{itemize}

\begin{tcolorbox}[breakable,colback=teal!4,colframe=teal!35!black,title={Before you move on}]
What does \mw{जंगळ} mean? (\textquotedblleft{}A forest\textquotedblright{}.) Say it once more.
\end{tcolorbox}
\section[\texorpdfstring{\mw{नदी} (nadī)}{नदी (nadī)}]{\mw{नदी} (nadī) — a river}

\label{lesson:MW-C68-nadi}



\begin{quote}
Before the new one: say the Marwadi for a hill, then the Marwadi for a forest.
\end{quote}

\subsection*{You'll want to know: \mw{नदी}}
\textbf{\mw{नदी}} — \emph{nadī} — \textquotedblleft{}a river\textquotedblright{}.

\begin{cousinweb}
\textbf{\mw{नदी}} (\emph{nadī}) is Sanskrit \emph{nadī}, \textquotedblleft{}river\textquotedblright{}, taken over whole as a learned loan, a tatsama (CDIAL 6943).
\end{cousinweb}

\begin{grammarlens}[title={What you've built}]
One more word for this chapter.
\end{grammarlens}

\subsection*{Your turn}
\begin{itemize}
\raggedright
  \item \emph{Say it:} \emph{nadī}
  \item \emph{Say it:} \emph{nadī}, once more
  \item \emph{Say it:} say \emph{nadī}
  \item \emph{Recall:} say the Marwadi for a hill, then the Marwadi for a forest, then say \emph{nadī} again
\end{itemize}

\begin{tcolorbox}[breakable,colback=teal!4,colframe=teal!35!black,title={Before you move on}]
What does \mw{नदी} mean? (\textquotedblleft{}A river\textquotedblright{}.) Say it once more.
\end{tcolorbox}
\section[\texorpdfstring{\mw{भाटो} (bhāṭo)}{भाटो (bhāṭo)}]{\mw{भाटो} (bhāṭo) — a stone}

\label{lesson:MW-C68-bhato}



\begin{quote}
Before the new one: say the Marwadi for a forest, then the Marwadi for a river.
\end{quote}

\subsection*{You'll want to know: \mw{भाटो}}
\textbf{\mw{भाटो}} — \emph{bhāṭo} — \textquotedblleft{}a stone\textquotedblright{}.

\begin{cousinweb}
The origin of \textbf{\mw{भाटो}} (\emph{bhāṭo}) is not settled. CDIAL, searched by form and by meaning, has no entry for it.
\end{cousinweb}

\begin{grammarlens}[title={What you've built}]
One more word for this chapter.
\end{grammarlens}

\subsection*{Your turn}
\begin{itemize}
\raggedright
  \item \emph{Say it:} \emph{bhāṭo}
  \item \emph{Say it:} \emph{bhāṭo}, once more
  \item \emph{Say it:} say \emph{bhāṭo}
  \item \emph{Recall:} say the Marwadi for a forest, then the Marwadi for a river, then say \emph{bhāṭo} again
\end{itemize}

\begin{tcolorbox}[breakable,colback=teal!4,colframe=teal!35!black,title={Before you move on}]
What does \mw{भाटो} mean? (\textquotedblleft{}A stone\textquotedblright{}.) Say it once more.
\end{tcolorbox}
\section[\texorpdfstring{\mw{रेत} (ret)}{रेत (ret)}]{\mw{रेत} (ret) — sand}

\label{lesson:MW-C68-ret}



\begin{quote}
Before the new one: say the Marwadi for a river, then the Marwadi for a stone.
\end{quote}

\subsection*{You'll want to know: \mw{रेत}}
\textbf{\mw{रेत}} — \emph{ret} — \textquotedblleft{}sand\textquotedblright{}.

\begin{cousinweb}
\textbf{\mw{रेत}} (\emph{ret}) is a tadbhava, a word inherited from Sanskrit and reshaped by sound change. It comes from \emph{retra}, \textquotedblleft{}perfumed powder\textquotedblright{}. CDIAL lists Old Marwari \emph{reta} under it (CDIAL 10816). CDIAL thinks a Munda source for the Sanskrit word scarcely likely.
\end{cousinweb}

\begin{grammarlens}[title={What you've built}]
One more word for this chapter.
\end{grammarlens}

\subsection*{Your turn}
\begin{itemize}
\raggedright
  \item \emph{Say it:} \emph{ret}
  \item \emph{Say it:} \emph{ret}, once more
  \item \emph{Say it:} say \emph{ret}
  \item \emph{Recall:} say the Marwadi for a river, then the Marwadi for a stone, then say \emph{ret} again
\end{itemize}

\begin{tcolorbox}[breakable,colback=teal!4,colframe=teal!35!black,title={Before you move on}]
What does \mw{रेत} mean? (\textquotedblleft{}Sand\textquotedblright{}.) Say it once more.
\end{tcolorbox}
\section[\texorpdfstring{\mw{खेत} (khet)}{खेत (khet)}]{\mw{खेत} (khet) — a field}

\label{lesson:MW-C68-khet}



\begin{quote}
Before the new one: say the Marwadi for a stone, then the Marwadi for sand.
\end{quote}

\subsection*{You'll want to know: \mw{खेत}}
\textbf{\mw{खेत}} — \emph{khet} — \textquotedblleft{}a field\textquotedblright{}.

\begin{cousinweb}
\textbf{\mw{खेत}} (\emph{khet}) is a tadbhava, a word inherited from Sanskrit and reshaped by sound change. It comes from \emph{kṣetra}, \textquotedblleft{}land\textquotedblright{}. CDIAL lists Marwari \emph{khet} under it (CDIAL 3735).
\end{cousinweb}

\begin{grammarlens}[title={What you've built}]
One more word for this chapter.
\end{grammarlens}

\subsection*{Your turn}
\begin{itemize}
\raggedright
  \item \emph{Say it:} \emph{khet}
  \item \emph{Say it:} \emph{khet}, once more
  \item \emph{Say it:} say \emph{khet}
  \item \emph{Recall:} say the Marwadi for a stone, then the Marwadi for sand, then say \emph{khet} again
\end{itemize}

\begin{tcolorbox}[breakable,colback=teal!4,colframe=teal!35!black,title={Before you move on}]
What does \mw{खेत} mean? (\textquotedblleft{}A field\textquotedblright{}.) Say it once more.
\end{tcolorbox}
